
\chapter{Introduction}

\noindent
IN THIS CHAPTER we begin our study of differential equations.

\medskip\noindent
 SECTION~1.1 presents
examples of applications that lead to differential equations.

\medskip\noindent
SECTION~1.2 introduces   basic concepts and
definitions concerning differential equations.

\medskip\noindent
 SECTION~1.3
presents a geometric method for dealing with differential equations
that has been known for a very long time, but has become particularly
useful and important with the proliferation of readily available
differential equations software.

\newpage

\section{Applications Leading to Differential Equations}

In order to apply mathematical methods to a physical or ``real life''
problem, we must  formulate the problem in mathematical
terms; that is, we must construct a \emph{\color{blue}mathematical
model} for the
problem. Many physical problems concern relationships between changing
quantities. Since rates of change are represented mathematically by
derivatives, mathematical models often involve equations relating an
unknown function and one or more of its derivatives. Such equations
are  \emph{\color{blue}differential equations}. They are the subject of this
book.

Much of calculus is devoted to learning mathematical techniques that
are applied in later courses in mathematics and the sciences;     you
wouldn't have time to learn much calculus if you insisted on seeing
a specific application of every topic covered in the course.
Similarly, much of this book is devoted to methods that can be applied
in later courses. Only a relatively small part of the book is devoted
to the derivation of specific differential equations from mathematical
models, or relating the differential equations that we study to
specific applications. In this section we mention a few such
applications.

The mathematical model for an applied problem is almost always simpler
than the actual situation being studied, since simplifying assumptions
are usually required to obtain a mathematical problem that can be
solved. For example, in modeling the motion of a falling object, we
might neglect air resistance and the gravitational pull of celestial
bodies other than Earth, or in  modeling  population growth we
might assume that the population grows continuously rather than in
discrete steps.

A good mathematical model has two important properties:

\begin{itemize}
\item It's sufficiently simple so that the mathematical problem
can be solved.

\item It represents the actual situation sufficiently well so that the
solution to the mathematical problem predicts the outcome of the real
problem to within a useful degree of accuracy. If results predicted by
the model don't agree with physical observations, the
underlying
assumptions of the model must be revised until satisfactory agreement
is obtained.
\end{itemize}

We'll  now give  examples of mathematical models involving
differential equations. We'll return to these problems at the
appropriate times, as we learn how to solve the various types of
differential equations that occur in the models.

All the examples in this section deal with functions of time, which we
 denote by $t$. If $y$ is a function of $t$,  $y'$
denotes the derivative of $y$ with respect to $t$;   thus,
\[
y'=\frac{dy}{dt}.
\]

\boxit{Population Growth and Decay}

\noindent
Although the number of members of a population (people in a given
country, bacteria in a laboratory culture, wildflowers in a forest,
etc.)
at any given time $t$ is necessarily an integer, models that use
differential equations to describe the growth and decay of populations
usually rest on the simplifying assumption that the number of members of
the population can be regarded as a differentiable function $P=P(t)$.
In most models it is assumed that the differential equation takes the
form
\begin{equation} \label{eq:1.1.1}
P'=a(P)P,
\end{equation}
where $a$ is a continuous function of $P$ that represents the rate of
change of population per unit time per individual.
In the
\href{http://en.wikipedia.org/wiki/Thomas_Robert_Malthus}
{\color{blue}\itshape Malthusian model},
 it is assumed that
$a(P)$ is a constant, so \eqref{eq:1.1.1} becomes
\begin{equation} \label{eq:1.1.2}
P'=aP.
\end{equation}
{\color{blue}
(When you see a name in blue italics, just click on it for
information about the person.)}
This model assumes that the numbers of births and deaths per unit time
are both proportional to the population. The constants of
proportionality are  the \emph{\color{blue}birth rate} (births per unit
time per individual) and the \emph{\color{blue}death rate} (deaths per unit time
per individual);     $a$ is the birth rate minus the death rate.
You  learned in calculus that if $c$ is any constant then
\begin{equation} \label{eq:1.1.3}
P=ce^{at}
\end{equation}
satisfies \eqref{eq:1.1.2}, so \eqref{eq:1.1.2} has infinitely
many solutions. To select the solution of the specific problem that
we're considering, we must know the population $P_0$ at an initial
time, say $t=0$. Setting $t=0$ in \eqref{eq:1.1.3} yields
$c=P(0)=P_0$, so the applicable solution is
\[
P(t)=P_0e^{at}.
\]
This implies that
\[
\lim_{t\to\infty}P(t)=\left\{\begin{array}{cl}\infty&\mbox{ if
}a>0,\\ 0&\mbox{ if }a<0;    \end{array}\right.
\]
that is, the population approaches infinity if the birth rate exceeds
the death rate, or zero if the death rate exceeds the birth rate.

To see the limitations of the Malthusian model, suppose  we're
modeling the  population of a country, starting from a time
$t=0$ when the birth rate exceeds the death rate (so $a>0$), and
the country's resources in terms of space, food supply, and other
necessities of life can support the existing population. Then the
prediction $P=P_0e^{at}$ may be reasonably accurate as long as
it remains within limits that the country's resources can support.
However, the model must inevitably lose validity when the prediction
exceeds these limits. (If nothing else, eventually there won't be
enough space for the predicted population!)

This flaw in the Malthusian model suggests the need for a model that
accounts for limitations of space and resources that tend to oppose
the rate of population growth as the population increases.
 Perhaps the most famous model of this kind is the
\href{http://www-history.mcs.st-and.ac.uk/Mathematicians/Verhulst.html}
{\color{blue}\itshape Verhulst model}, where \eqref{eq:1.1.2} is
replaced by
\begin{equation} \label{eq:1.1.4}
P'=aP(1-\alpha P),
\end{equation}
where $\alpha$ is a positive constant.
As long as $P$ is small
compared to $1/\alpha$, the ratio $P'/P$  is approximately equal to $a$.
Therefore the growth is approximately exponential;     however, as $P$
increases, the ratio $P'/P$  decreases as opposing factors
become significant.

Equation~\eqref{eq:1.1.4} is  the \emph{\color{blue}logistic equation}. You
will learn how to solve it in Section~1.2. (See
Exercise~2.2.~\hspace*{-3pt}\ref{exer:2.2.28}.) The solution is
\[
P=\frac{P_0}{\alpha P_0+(1-\alpha P_0)e^{-at}},
\]
where $P_0=P(0)>0$. Therefore
$\lim_{t\to\infty}P(t)=1/\alpha$, independent of $P_0$.

Figure~\ref{figure:1.1.1} shows typical graphs of $P$ versus $t$ for
various values of  $P_0$.

\begin{figure}[h]
  \centering
  \includegraphics[%bb=-78 148 689 643,
  width=5.67in,height=3.66in,keepaspectratio]{fig010101}
\color{blue}
  \caption{Solutions of the logistic equation}
  \label{figure:1.1.1}
\end{figure}


\boxit{Newton's Law of Cooling}

\noindent
According to
\href{http://www-history.mcs.st-and.ac.uk/Mathematicians/Newton.html}
{\color{blue}\itshape Newton's law of cooling},  the temperature of a
body changes at a rate proportional to the difference between the
temperature of the body and the temperature of the surrounding medium.
 Thus, if  $T_m$ is the temperature of the
medium and
$T=T(t)$ is the temperature of the body at time $t$, then
\begin{equation} \label{eq:1.1.5}
T' = -k(T-T_m),
\end{equation}
where $k$ is a positive constant and the  minus sign indicates;   that
the temperature of the body increases with time if it's less than the
temperature of the medium, or decreases if it's  greater. We'll see
in Section~4.2 that if
$T_m$ is constant then the solution of \eqref{eq:1.1.5} is
\begin{equation} \label{eq:1.1.6}
T=T_m+(T_0-T_m)e^{-kt},
\end{equation}
where $T_0$ is the temperature of the body when $t=0$.
Therefore $\lim_{t\to\infty}T(t)=T_m$, independent of $T_0$.
(Common sense suggests this. Why?)

Figure~\ref{figure:1.1.2} shows typical graphs of $T$ versus $t$ for
various values of  $T_0$.

\begin{figure}[htp]
  \centering
  \includegraphics[%bb=-78 148 689 643,
  width=5.67in,height=3.66in,keepaspectratio]{fig010102}
\color{blue}
  \caption{ Temperature according to Newton's Law of Cooling}
  \label{figure:1.1.2}
\end{figure}

Assuming that the medium remains at constant temperature seems
reasonable if we're considering a cup of coffee cooling in a room, but
 not if we're cooling a huge cauldron of molten
metal in the same room. The difference between the two situations is
that the heat lost by the coffee isn't likely to raise the
temperature of the room appreciably, but the heat lost by the
cooling metal is. In this second situation we must use a model that
accounts for the heat exchanged between the object and the medium. Let
$T=T(t)$ and $T_m=T_m(t)$ be the temperatures of the object and the
medium respectively, and let $T_0$ and $T_{m0}$ be their initial
values. Again, we assume that $T$ and $T_m$ are related by
\eqref{eq:1.1.5}. We also assume that the change in heat of
the object as its temperature changes from $T_0$ to $T$ is $a(T-T_0)$
and  the change in heat of the medium as its temperature changes
from $T_{m0}$ to $T_m$ is $a_m(T_m-T_{m0})$, where $a$ and $a_m$ are
positive constants depending upon the masses and thermal properties of
the object and medium respectively. If we assume that the total heat
of the in the object and the medium remains constant
(that is, energy is conserved), then
\[
a(T-T_0)+a_m(T_m-T_{m0})=0.
\]
Solving this for  $T_m$  and substituting the result into
\eqref{eq:1.1.6} yields the differential equation
\[
T'=-k\left(1+\frac{a}{a_m}\right)T
+k\left(T_{m0}+\frac{a}{a_m}T_0\right)
\]
for the temperature of the object.  After learning to solve linear
first order  equations, you'll be able to show
(Exercise~4.2.~\hspace*{-3pt}\ref{exer:4.2.17}) that
\[
T=\frac{aT_0+a_mT_{m0}}{a+a_m}+\frac{a_m(T_0-T_{m0})}{a+a_m}e^{-k(1+a/a_m)t}.
\]

\boxit{Glucose Absorption by the Body}

\noindent
Glucose is absorbed by
the body at a rate proportional to the amount of glucose present in
the bloodstream. Let $\lambda$ denote the (positive) constant of
proportionality. Suppose   there are $G_0$ units of glucose in
the bloodstream when $t=0$, and let $G=G(t)$ be the number of units in
the bloodstream at time $t>0$. Then, since the glucose being absorbed
by the body is leaving the bloodstream, $G$ satisfies the equation
\begin{equation} \label{eq:1.1.7}
G'=-\lambda G.
\end{equation}
From  calculus you know that if $c$ is any constant then
\begin{equation} \label{eq:1.1.8}
G=ce^{-\lambda t}
\end{equation}
satisfies \eqref{eq:1.1.7}, so \eqref{eq:1.1.7} has infinitely
many solutions.
 Setting $t=0$ in \eqref{eq:1.1.8} and requiring that
$G(0)=G_0$ yields $c=G_0$, so
\[
G(t)=G_0e^{-\lambda t}.
\]

Now let's complicate matters by injecting glucose intravenously
at a constant rate of $r$ units of glucose per unit of time.
Then the rate of change of the amount of glucose  in the bloodstream
per unit time is
\begin{equation} \label{eq:1.1.9}
G'=-\lambda G+r,
\end{equation}
where the first term on the right is due to the absorption of the
glucose by the body and the second term is due to the injection.
 After you've studied Section~2.1,
you'll be able to show (Exercise~2.1.~\hspace*{-3pt}\ref{exer:2.1.43}) that the solution
of
\eqref{eq:1.1.9} that satisfies $G(0)=G_0$ is
\[
G=\frac{r}{\lambda}+\left(G_0-\frac{r}{\lambda}\right)e^{-\lambda t}.
\]
Graphs of  this function are similar to those in
Figure~\ref{figure:1.1.2}.
(Why?)

\boxit{Spread of Epidemics}

\noindent
One model for the spread of epidemics assumes that the number of
people infected changes at a rate proportional to the product of the
number of people already infected and the number of people who are
susceptible, but not yet infected. Therefore, if $S$ denotes the
total population of susceptible people and $I=I(t)$ denotes the
number
of infected people at time $t$, then $S-I$ is the number of people
who are susceptible, but not yet infected. Thus,
\[
I'=rI(S-I),
\]
where $r$ is a positive constant. Assuming that $I(0)=I_0$,
the solution of this equation is
\[
I=\frac{SI_0}{I_0+(S-I_0)e^{-rSt}}
\]
(Exercise~2.2.~\hspace*{-3pt}\ref{exer:2.2.29}).
 Graphs of this function are similar to those in
Figure~\ref{figure:1.1.1}.
(Why?)
Since $\lim_{t\to\infty}I(t)=S$, this model predicts that all the
susceptible people eventually become infected.

\boxit{Newton's Second Law of Motion}

\noindent

According to
\href{http://www-history.mcs.st-and.ac.uk/Mathematicians/Newton.html}
{\color{blue}\itshape Newton's second law of motion},  the
instantaneous acceleration
$a$ of an object with constant mass $m$ is related to the force $F$
acting on the object by the equation $F=ma$. For simplicity, let's
assume that $m=1$ and the motion of the object is along a vertical
line. Let $y$ be the displacement of the object from some reference
point on Earth's surface, measured positive upward. In many
applications, there are three kinds of forces that may act on the
object:

\begin{alist}
\item % (a)
A force such as gravity that depends only on the position $y$,
which we write as $-p(y)$, where $p(y)>0$ if $y\ge0$.

\item % (b)
A force such as atmospheric resistance that depends on
the position and velocity of the object, which we write as
$-q(y,y')y'$, where $q$ is a nonnegative function and we've
put $y'$ ``outside'' to indicate that the resistive force is
always in the direction opposite to the velocity.
\item % (c)
A force $f=f(t)$, exerted from an external source (such as a towline
from a helicopter) that depends only on $t$.
\end{alist}

In this case, Newton's second law implies that
\[
y''=-q(y,y')y'-p(y)+f(t),
\]
which is usually rewritten as
\[
y''+q(y,y')y'+p(y)=f(t).
\]
Since the  second (and no higher) order derivative of $y$ occurs in
this equation, we say that it is a \emph{\color{blue}second order differential
equation}.

\boxit{Interacting Species: Competition}

\noindent
Let $P=P(t)$ and $Q=Q(t)$ be the populations of two species at time
$t$, and assume that each population would grow exponentially if the
other didn't exist; that is, in the absence of competition we would
have
\begin{equation} \label{eq:1.1.10}
P'=aP \mbox{\quad and \quad}Q'=bQ,
\end{equation}
where $a$ and $b$ are positive constants. One way to model the effect
of competition is to assume that the growth rate per individual of
each population is reduced by an amount proportional to the other
population, so \eqref{eq:1.1.10} is replaced by
\begin{eqnarray*}
P'&=&\phantom{-}aP-\alpha Q\\
Q'&=&-\beta P+bQ,
\end{eqnarray*}
where $\alpha$ and $\beta$ are positive constants. (Since negative
population doesn't make sense, this system works only while $P$ and
$Q$ are both positive.) Now suppose   $P(0)=P_0>0$ and
$Q(0)=Q_0>0$. It can be shown (Exercise~10.4.~\hspace*{-3pt}\ref{exer:10.4.42})
that there's a  positive constant $\rho$ such that if
$(P_0,Q_0)$ is above the line $L$ through the origin with slope $\rho$,
then the species with population $P$ becomes extinct in finite time,
but if $(P_0,Q_0)$ is below $L$,   the species with population
$Q$ becomes extinct in finite time. Figure~\ref{figure:1.1.3} illustrates
this. The curves shown there are given parametrically by $P=P(t),
Q=Q(t),\ t>0$.
 The arrows indicate direction along the curves with
increasing $t$.

\begin{figure}[h]
  \centering
  \includegraphics[%bb=-78 148 689 643,
  width=5.67in,height=3.66in,keepaspectratio]{fig010103}
\color{blue}
  \caption{ Populations of competing species}
  \label{figure:1.1.3}
\end{figure}

\section{Basic Concepts}

A \emph{\color{blue}differential equation} is an equation that contains one or more
derivatives of an unknown function. The \emph{\color{blue}order} of a
differential equation is the order of the highest
derivative that it contains. A differential equation is an
\emph{\color{blue}ordinary differential equation} if it involves an unknown
function of only one variable, or a \emph{\color{blue}partial differential
equation} if it involves partial derivatives of a function of more
than one variable. For now we'll consider only ordinary differential
equations, and we'll just call them \emph{\color{blue}differential equations}.

Throughout this text, all variables and constants are real unless it's
stated otherwise. We'll usually use $x$ for the independent variable
unless the independent variable is time; then we'll use $t$.

The simplest differential equations are first order equations of the
form
\[
\frac{dy}{dx}=f(x) \mbox{\quad or, equivalently, \quad}y'=f(x),
\]
where $f$  is a known function of $x$. We already know from calculus
how to find functions that satisfy this kind of equation. For example,
if
\[
y'=x^3,
\]
then
\[
y=\int x^3\, dx=\frac{x^4}{4}+c,
\]
 where $c$ is an arbitrary constant.  If $n>1$
we can find functions $y$ that satisfy
equations of the form
\begin{equation} \label{eq:1.2.1}
y^{(n)}=f(x)
\end{equation}
by repeated integration. Again, this is a calculus problem.

Except for illustrative purposes in this section, there's no need to
consider differential equations like \eqref{eq:1.2.1}.We'll
usually consider differential equations that can be written as
\begin{equation} \label{eq:1.2.2}
y^{(n)}=f(x,y,y', \dots,y^{(n-1)}),
\end{equation}
where at least one of the functions $y$, $y'$, \dots, $y^{(n-1)}$ actually
appears on the right. Here are some examples:
\begin{align*}
\frac{dy}{dx}-x^2&=0&&\text{(first order)},  \\
\frac{dy}{dx}+2xy^2&=-2&&\text{(first order)},   \\
\frac{d^2y}{dx^2}+2\frac{dy}{dx}+y&=2x&&\text{(second order)},\\
xy'''+y^2&=\sin x  &&\text{(third order)},\\
y^{(n)}+xy'+3y&=x&&\text{($n$-th order)}.
\end{align*}
Although none of these equations is  written as in
\eqref{eq:1.2.2}, all of them \emph{\color{blue}can} be written in this form:
\begin{align*}
y'&=x^2,  \\
y'&=-2-2xy^2,   \\
y''&=2x-2y'-y,  \\
y'''&=\frac{\sin x-y^2}{x},\\
y^{(n)}&=x-xy'-3y.
\end{align*}

\boxit{Solutions of Differential Equations}

\noindent
A \emph{\color{blue}solution} of a differential equation is a function that
satisfies the differential equation on some open interval;   thus, $y$
is a solution of \eqref{eq:1.2.2} if $y$ is $n$ times differentiable and
\[
y^{(n)}(x)=f(x,y(x),y'(x), \dots,y^{(n-1)}(x))
\]
for all $x$ in some open interval $(a,b)$. In this case, we also say
that $y$ \emph{\color{blue}is a solution of $\eqref{eq:1.2.2}$ on} $(a,b)$. Functions
that satisfy a differential equation at isolated points are not
interesting. For example, $y=x^2$ satisfies
\[
xy'+x^2=3x
\]
if and only if $x=0$ or $x=1$, but it's not a solution of this
differential equation because it does not satisfy the equation on an
open interval.

The graph of a solution of a differential equation is  a {\color{blue}\it
solution curve\/}. More generally, a curve $C$ is said to be an {\color{blue}\it
integral curve\/} of a differential equation if every function
$y=y(x)$ whose graph is a segment of $C$ is a solution of the
differential equation. Thus, any solution curve of a differential
equation is an integral curve, but an integral curve need not be a
solution curve.

\begin{example}
\label{example:1.2.1} \rm
If $a$ is any positive constant,  the circle
\begin{equation} \label{eq:1.2.3}
x^2+y^2=a^2
\end{equation}
is an integral curve of
\begin{equation} \label{eq:1.2.4}
y'=-\frac{x}{y}.
\end{equation}
To see this, note that the only functions whose graphs are segments of
\eqref{eq:1.2.3} are
\[
y_1=\sqrt{a^2-x^2}\mbox{\quad and \quad} y_2=-\sqrt{a^2-x^2}.
\]
We leave it to you to verify that these functions both satisfy
\eqref{eq:1.2.4} on the open interval $(-a,a)$. However, \eqref{eq:1.2.3} is
not a solution
curve of \eqref{eq:1.2.4}, since it's not the graph of a function.
\end{example}

\begin{example}\label{example:1.2.2} \rm
Verify that
\begin{equation} \label{eq:1.2.5}
y=\frac{x^2}{3}+\frac{1}{x}
\end{equation}
is a solution of
\begin{equation} \label{eq:1.2.6}
xy'+y=x^2
\end{equation}
 on $(0,\infty)$ and on $(-\infty,0)$.
\end{example}

\solution
Substituting \eqref{eq:1.2.5} and
\[
y'=\frac{2x}{3} - \frac{1}{x^2}
\]
into \eqref{eq:1.2.6}  yields
\[
xy'(x)+y(x)=x \left(\frac{2x}{3} - \frac{1}{x^2}\right)+
\left(\frac{x^2}{3}+\frac{1}{x}\right)=x^2
\]
for all $x\ne0$. Therefore $y$ is a solution of \eqref{eq:1.2.6}
on $(-\infty,0)$ and $(0,\infty)$.
 However, $y$ isn't  a solution of the differential
equation on any open interval that contains $x=0$, since  $y$ is
not defined at  $x=0$.

Figure~\ref{figure:1.2.1} shows the graph
of \eqref{eq:1.2.5}. The
part
of the graph of \eqref{eq:1.2.5} on $(0,\infty)$ is a solution curve of
\eqref{eq:1.2.6}, as is the part of the graph on $(-\infty,0)$.

\begin{example}\label{example:1.2.3} \rm
 Show that if $c_1$ and $c_2$ are constants then
\begin{equation} \label{eq:1.2.7}
y=(c_1+c_2x)e^{-x}+2x-4
\end{equation}
 is a solution of
\begin{equation} \label{eq:1.2.8}
y''+2y'+y=2x
\end{equation}
 on $(-\infty,\infty)$.
\end{example}

\begin{figure}[thp]
  \centering
  \includegraphics[%bb=-78 148 689 643,
  width=5.67in,height=3.66in,keepaspectratio]{fig010201}
\color{blue}
  \caption{$y=\dst{\frac{x^{2}}{3}+\frac{1}{x}}$}
  \label{figure:1.2.1}
\end{figure}

\solution Differentiating \eqref{eq:1.2.7} twice yields
\[
y'=-(c_1+c_2x)e^{-x}+c_2e^{-x}+2
\]
 and
\[
y''=(c_1+c_2x)e^{-x}-2c_2e^{-x},
\]
so
\begin{eqnarray*}
y''+2y'+y&=&(c_1+c_2x)e^{-x}-2c_2e^{-x}\\
&&+2\left[-(c_1+c_2x)e^{-x}+c_2e^{-x}+2\right]\\
&&+(c_1+c_2x)e^{-x}+2x-4\\
&=&(1-2+1)(c_1+c_2x)e^{-x}+(-2+2)c_2e^{-x}\\ &&+4+2x-4=2x
\end{eqnarray*}
for all values of $x$.
Therefore $y$ is a solution of \eqref{eq:1.2.8} on $(-\infty,\infty)$.

\begin{example}\label{example:1.2.4} \rm
Find all solutions of
\begin{equation} \label{eq:1.2.9}
y^{(n)}=e^{2x}.
\end{equation}
\end{example}

\solution Integrating \eqref{eq:1.2.9} yields
\[
y^{(n-1)}=\frac{e^{2x}}{2}+k_1,
\]
 where $k_1$ is a constant. If $n\ge2$,
integrating again yields
\[
y^{(n-2)}=\frac{e^{2x}}{4}+k_1x+k_2.
\]
If $n\ge3$, repeatedly integrating yields
\begin{equation} \label{eq:1.2.10}
y=\frac{e^{2x}}{2^n}+k_1\frac{x^{n-1}}{(n-1)!}+k_2\frac{x^{n-2}}{(n-2)!}+\cdots+k_n,
\end{equation}
 where $k_1$, $k_2$, \dots, $k_n$ are constants.
This shows that every solution of \eqref{eq:1.2.9} has the form
\eqref{eq:1.2.10}
for some choice of the constants $k_1$, $k_2$, \dots, $k_n$.
On the other hand, differentiating \eqref{eq:1.2.10}  $n$ times shows
that if
$k_1$, $k_2$, \dots, $k_n$ are arbitrary constants, then the function $y$
in
\eqref{eq:1.2.10} satisfies \eqref{eq:1.2.9}.

Since the constants $k_1$, $k_2$, \dots, $k_n$  in \eqref{eq:1.2.10}
are arbitrary, so are the constants
\[\frac{k_1}{(n-1)!},\, \frac{k_2}{(n-2)!},\, \cdots, \, k_n.\]
 Therefore Example~\ref{example:1.2.4} actually shows that all
solutions of \eqref{eq:1.2.9}  can be written as
\[
y=\frac{e^{2x}}{2^n}+c_1+c_2x+\cdots+c_nx^{n-1},
\]
 where we  renamed the arbitrary constants in
\eqref{eq:1.2.10} to obtain a simpler formula. As a
general rule, arbitrary constants appearing in solutions  of differential
equations should be simplified if possible. You'll see examples
of this throughout the text.

\boxit{Initial Value Problems}

\noindent
In Example~\ref{example:1.2.4} we saw that the differential equation
$y^{(n)}=e^{2x}$ has an infinite family of solutions that depend upon
the $n$ arbitrary constants $c_1$, $c_2$, \dots, $c_n$. In the absence of
additional conditions, there's no reason to prefer one solution of a
differential equation over another. However, we'll often be
interested in finding a solution of a differential equation that
satisfies one or more specific conditions. The next example
illustrates this.

\begin{example}\label{example:1.2.5} \rm
 Find a solution of
\[
y'=x^3
\]
 such that $y(1)=2$.
\end{example}

\solution  At the beginning of this section we saw
that the  solutions of $y'=x^3$ are
\[
y=\frac{x^4}{4}+c.
\]
 To determine a value of $c$ such that $y(1)=2$,
we set $x=1$ and $y=2$ here to obtain
\[
2=y(1)=\frac{1}{4}+c,\mbox{\quad so \quad}c=\frac{7}{4}.
\]
 Therefore the required solution is
\[
y=\frac{x^4+7}{4}.
\]

Figure~\ref{figure:1.2.2} shows
the graph  of this solution.
Note that imposing the  condition $y(1)=2$ is equivalent to requiring
 the graph of $y$ to pass through the point $(1,2)$.

 We can rewrite the problem considered in Example~\ref{example:1.2.5}
more briefly as
\[
y'=x^3,\quad y(1)=2.
\]
 We call this an \emph{\color{blue}initial value problem}.
The requirement $y(1)=2$ is an \emph{\color{blue}initial condition}.
 Initial value problems can also be
posed for higher order differential equations.  For example,
\begin{equation} \label{eq:1.2.11}
y'' - 2y'+3y=e^x, \quad y(0)=1, \quad y'(0)=2
\end{equation}
is  an initial value problem for a second order differential
equation where $y$ and $y'$ are required to have specified values at
 $x=0$. In general, an initial value
problem for an $n$-th order differential equation requires $y$ and its
first $n-1$ derivatives to have specified values at some point $x_0$.
These requirements are the \emph{\color{blue}initial conditions}.

\begin{figure}[H]
  \centering
  \includegraphics[%bb=-78 148 689 643,
  width=5.67in,height=3.66in,keepaspectratio]{fig010202}
\color{blue}
  \caption{$y=\dst{\frac{x^{2}+7}{4}}$ }
  \label{figure:1.2.2}
\end{figure}

We'll denote an initial value problem for a differential equation by
writing the initial conditions after the equation, as in
\eqref{eq:1.2.11}. For example, we would write an initial value problem
for \eqref{eq:1.2.2} as
\begin{equation} \label{eq:1.2.12}
y^{(n)}=f(x,y,y', \dots,y^{(n-1)}),\, y(x_0)=k_0,\,
y'(x_0)=k_1,\, \dots,\, y^{(n-1)}=k_{n-1}.
\end{equation}
Consistent with our earlier definition of a solution of the
differential equation in \eqref{eq:1.2.12}, we say that $y$ is a solution
of the initial value problem \eqref{eq:1.2.12}  if $y$ is $n$ times
differentiable and
\[
y^{(n)}(x)=f(x,y(x),y'(x), \dots,y^{(n-1)}(x))
\]
for all $x$ in some open interval $(a,b)$  that contains $x_0$,
and  $y$ satisfies the initial conditions in \eqref{eq:1.2.12}. The
largest open interval that contains $x_0$ on which $y$ is defined and
satisfies the differential equation is  the \emph{\color{blue}interval of
validity} of $y$.

\begin{example}\label{example:1.2.6} \rm
In Example~\ref{example:1.2.5} we saw that
\begin{equation} \label{eq:1.2.13}
y=\frac{x^4+7}{4}
\end{equation}
is a solution of the initial value problem
\[
y'=x^3,\quad y(1)=2.
\]
Since the function in \eqref{eq:1.2.13} is defined for all $x$, the
interval of validity of this solution is $(-\infty,\infty)$.
\end{example}

\begin{example}\label{example:1.2.7} \rm
In Example~\ref{example:1.2.2} we verified that
\begin{equation} \label{eq:1.2.14}
y=\frac{x^2}{3}+\frac{1}{x}
\end{equation}
is a solution of
\[
xy'+y=x^2
\]
on $(0,\infty)$ and on $(-\infty,0)$. By evaluating \eqref{eq:1.2.14} at
$x=\pm1$, you can see that \eqref{eq:1.2.14} is a solution of the initial
value problems
\begin{equation} \label{eq:1.2.15}
xy'+y=x^2,\quad y(1)=\frac{4}{3}
\end{equation}
and
\begin{equation} \label{eq:1.2.16}
xy'+y=x^2,\quad y(-1)=-\frac{2}{3}.
\end{equation}
The interval of validity of \eqref{eq:1.2.14} as a solution of
\eqref{eq:1.2.15} is $(0,\infty)$, since this is the largest interval
that contains $x_0=1$ on which \eqref{eq:1.2.14} is defined. Similarly, the
interval of validity of \eqref{eq:1.2.14} as a solution of \eqref{eq:1.2.16}
is $(-\infty,0)$, since this is the largest interval that contains
$x_0=-1$ on which \eqref{eq:1.2.14} is defined.
\end{example}

\boxit{Free Fall Under Constant Gravity}

\noindent
The term \emph{\color{blue}initial value problem}  originated in problems
of motion where the independent variable is $t$
(representing elapsed time), and the initial conditions
are the position and velocity of an object at the initial
(starting) time of an experiment.

\begin{example}\label{example:1.2.8} \rm
An object falls under the influence of gravity near Earth's
surface, where it can be assumed that the magnitude of the
acceleration due to gravity is a constant $g$.
\begin{alist}
\item % (a)
 Construct a mathematical model for the motion of the object in the
form of an initial value problem for a second order differential
equation, assuming that the altitude and velocity of the object at
time $t=0$ are known. Assume that gravity is the only force acting on
the object.
\item % (b)
 Solve the initial value problem derived in \part{a} to obtain the
altitude as a function of time.
\end{alist}
\end{example}

\solutionpart{a} Let $y(t)$ be the altitude of the object at time $t$.
 Since the acceleration of the object has
constant magnitude $g$ and is in the downward (negative) direction,
$y$ satisfies the second order equation
\[
y''=-g,
\]
where the prime now indicates differentiation with respect to $t$.
If $y_0$ and $v_0$ denote the altitude and velocity when
$t=0$, then $y$ is a solution of the initial value problem
\begin{equation} \label{eq:1.2.17}
y''=-g,\quad  y(0)=y_0,\quad y'(0)=v_0.
\end{equation}

\solutionpart{b}
Integrating \eqref{eq:1.2.17} twice yields
\begin{eqnarray*}
y'&=&-gt+c_1, \\
y&=&-\frac{gt^2}{2}+c_1t+c_2.
\end{eqnarray*}
Imposing the initial conditions $y(0)=y_0$ and $y'(0)=v_0$ in these
two equations shows that $c_1=v_0$ and $c_2=y_0$. Therefore the
solution of the initial value problem \eqref{eq:1.2.17} is
\[
y=- \frac{gt^2}{2}+v_0t+y_0.
\]

\exercises
\begin{exerciselist}
\item\label{exer:1.2.1}
Find the order of the   equation.

\begin{tabular}[t]{@{}p{168pt}@{}p{168pt}}
{\bf (a)}\   $\dst\frac{d^2y}{dx^2}+2
\frac{dy}{dx}\   \frac{d^3y}{dx^3}+x=0$
& {\bf (b)}\   $y''-3y'+2y=x^7$\\[9pt]
{\bf (c)}\   $y'-y^7=0$
& {\bf (d)}\   $y''y-(y')^2=2$
\end{tabular}

\item\label{exer:1.2.2}
Verify that the  function is a solution of the
differential equation on some interval, for any choice of
the arbitrary constants appearing in the function.

\begin{alist}
\item % (a)
 $y=ce^{2x};     \quad y'=2y$
\item % (b)
 $y=\dst\frac{x^2}{3}
+\frac{c}{x};   \quad xy'+y=x^2$

\item %(c)
$y=\dst\frac{1}{2}+ce^{-x^2};    \quad y'+2xy=x$

\item %(d)
$y=(1+ce^{-x^2/2});   (1-ce^{-x^2/2})^{-1}   \quad
2y'+x(y^2-1)=0$

\item %(e)
$y=\dst{\tan\left( \frac{x^3}{3}+c\right)};     \quad
y'=x^2(1+y^2)$

\item %(f)
$y=(c_1+c_2x)e^x+\sin x+x^2;    \quad
y''-2y'+y=-2 \cos x+x^2-4x+2$

\item %(g)
$y=c_1e^x+c_2x+\dst\frac{2}{x};     \quad
(1-x)y''+xy'- y=4(1-x-x^2)x^{-3}$

\item %(h)
$y=x^{-1/2}(c_1\sin x+c_2 \cos x)+4x+8$;     \newline
$x^2y''+xy'+\dst{\left(x^2-\frac{1}{4}\right)}y=4x^3+8x^2+3x-2$
\end{alist}

\item\label{exer:1.2.3}
Find all solutions of the  equation.

\begin{tabular}[t]{@{}p{168pt}@{}p{168pt}}
{\bf (a)}\quad  $y'=-x$
& {\bf (b)}\quad  $y'=-x \sin x$\\
{\bf (c)}\quad  $y'=x \ln x$
& {\bf (d)}\quad  $y''=x \cos x$\\
{\bf (e)}\quad  $y''=2xe^x$
& {\bf (f)}\quad  $y''=2x+\sin x+e^x$\\
{\bf (g)}\quad  $y'''=-\cos x$
& {\bf (h)}\quad  $y'''=-x^2+e^x$\\
{\bf (i)}\quad  $y'''=7e^{4x}$&
\end{tabular}

\item\label{exer:1.2.4}
Solve the  initial value problem.

\begin{alist}
\item % (a)
 $y'=-xe^x, \quad y(0)=1$
\item % (b)
 $\dst{y'=x \sin x^2, \quad y\left({\sqrt\frac{\pi}{2}}\right)=1}$

\item %(c)
$y'=\tan x, \quad y(\pi/4)=3$

\item %(d)
$y''=x^4, \quad y(2)=-1, \quad y'(2)=-1$

\item %(e)
$y''=xe^{2x}, \quad y(0)=7, \quad y'(0)=1$

\item %(f)
$y''=- x \sin x, \quad y(0)=1, \quad y'(0)=-3$

\item %(g)
$y'''=x^2e^x, \quad y(0)=1, \quad y'(0)=-2, \quad
y''(0)=3$

\item %(h)
$y'''=2+\sin 2x, \quad y(0)=1, \quad y'(0)=-6, \quad
y''(0)=3$

\item %(i)
$y'''=2x+1, \quad y(2)=1, \quad
y'(2)=-4, \quad y''(2)=7$
\end{alist}

\item\label{exer:1.2.5}
Verify that the  function is a solution of
the initial value problem.

\begin{alist}
\item %(a)
$y=x\cos x;     \quad y'=\cos x-y\tan x, \quad
y(\pi/4)=\dst\frac{\pi}{4\sqrt{2}}$

\item %(b)
$\dst{y=\frac{1+2\ln x}{x^2}+\frac{1}{2};     \quad
y'=\frac{x^2-2x^2y+2}{x^3},   \quad y(1)=\frac{3}{2}}$

\item %(c)
$y=\dst{\tan\left(\frac{x^2}{2}\right)};
\quad y'=x(1+y^2),   \quad y(0)=0$

\item %(d)
$\dst{y=\frac{2}{x-2};     \quad y'=\frac{-y(y+1)}{x}}, \quad
y(1)=-2$
\end{alist}

\item\label{exer:1.2.6}
Verify that the  function is a solution of
the initial value problem.

\begin{alist}
\item %(a)
$y=x^2(1+\ln x);     \quad y''=\dst\frac{3xy'-4y}{x^2}, \quad
y(e)=2e^2, \quad y'(e)=5e$

\item %(b)
$y=\dst\frac{x^2}{3}+x-1;     \quad y''=\dst\frac{x^2-xy'+y+1}{x^2}, \quad y(1)=\dst\frac{1}{3}, \quad y'(1)=\dst\frac{5}{3}$

\item %(c)
$y=(1+x^2)^{-1/2};     \quad
y''=\dst\frac{(x^2-1)y-x(x^2+1)y'}{(x^2+1)^2}, \quad y(0)=1,$

$y'(0)=0$

\item %(d)
$y=\dst\frac{x^2}{1-x};    \quad y''=\dst\frac{2(x+y)(xy'-y)}{x^3},
\quad y(1/2)=1/2, \quad y'(1/2)=3$
\end{alist}

\item\label{exer:1.2.7}
Suppose an object is launched from a point 320 feet
above the earth with an initial velocity of 128 ft/sec
upward, and  the only force acting on it thereafter is
gravity. Take $g=32$ ft/sec$^2$.

\begin{alist}
\item %(a)
Find the highest altitude attained by the object.

\item %(b)
Determine how long it takes for the object to fall to the
ground.
\end{alist}

\item\label{exer:1.2.8}
\setcounter{equation}{0}
Let $a$ be a nonzero real number.
\begin{alist}
\item % (a)
Verify that if $c$ is an arbitrary constant then
\[
y=(x-c)^a
\eqno{\rm (A)}
\]
is a solution of
\[
y'=ay^{(a-1)/a}
\eqno{\rm (B)}
\]
on $(c,\infty)$.
\item % (b)
Suppose $a<0$ or $a>1$. Can you think of a solution of
(B) that isn't of the form (A)?
\end{alist}

\item\label{exer:1.2.9}
Verify that
\[
y=
\left\{ \begin{array}{cl}
e^x-1,& x \ge 0, \\[6pt]
1-e^{-x},& x < 0, \end{array}\right.
\]
is a solution of
\[
y'=|y|+1
\]
on $(-\infty,\infty)$. \hint{Use the definition of derivative
at $x=0$.}

\item\label{exer:1.2.10}
\begin{alist}
\item %(a)
Verify that if $c$ is any real number then
\[
y=c^2+cx+2c+1
\eqno{\rm (A)}
\]
satisfies
\[
 y'=\frac{-(x+2)+\sqrt{x^2+4x+4y}}{2}
\eqno{\rm (B)}
\]
on some open interval. Identify the open interval.

\item %(b)
Verify that
\[
y_1=\frac{-x(x+4)}{4}
\]
also satisfies  (B) on some open interval, and
identify the open interval. (Note that $y_1$ can't be obtained
by selecting a value of $c$ in (A).)
\end{alist}

\end{exerciselist}

\section{Direction Fields for First Order Equations}

It's impossible to find explicit formulas for solutions of some
differential equations. Even if there are such  formulas, they may be
so complicated that they're useless. In this case we may resort to
graphical or numerical methods to get some idea of how the solutions
of the given equation behave.

 In Section~2.3 we'll take up
the question of existence of solutions of a first order equation
\begin{equation} \label{eq:1.3.1}
y'=f(x,y).
\end{equation}
In this section we'll simply assume that \eqref{eq:1.3.1} has
solutions and discuss a graphical method for approximating them.
In Chapter~3 we discuss numerical methods for obtaining
approximate solutions of \eqref{eq:1.3.1}.

Recall that a solution of \eqref{eq:1.3.1} is a function $y=y(x)$ such
that
\[
y'(x)=f(x,y(x))
\]
for all values of $x$ in some interval, and an integral curve is
either the graph of a solution or is
made up of  segments that are graphs of solutions.
 Therefore, not being able to
solve \eqref{eq:1.3.1} is equivalent to not knowing the equations of
integral curves of \eqref{eq:1.3.1}. However, it's easy to calculate the
slopes of these curves. To be specific, the slope of an integral curve
of \eqref{eq:1.3.1} through a given point $(x_0,y_0)$ is given by the
number $f(x_0,y_0)$. This is the basis of \emph{\color{blue}the method of direction fields}.

If $f$ is defined on a set $R$, we can construct a \emph{\color{blue}direction
field} for \eqref{eq:1.3.1}  in $R$ by drawing a
short
line segment through each point $(x,y)$ in $R$ with slope $f(x,y)$. Of
course, as a practical matter, we can't actually draw line segments
through \emph{\color{blue}every} point in $R$;     rather, we must select
a finite
set of points in $R$. For example, suppose   $f$ is defined on the
closed rectangular region
\[
R:\{a\le x\le b, c\le y\le d\}.
\]
Let
\[
a= x_0< x_1< \cdots< x_m=b
\]
be  equally spaced points in $[a,b]$ and
\[
c=y_0<y_1<\cdots<y_n=d
\]
be  equally spaced points in $[c,d]$.
 We say that the points
\[
 (x_i,y_j),\quad 0\le i\le m,\quad 0\le j\le n,
\]
form a \emph{\color{blue}rectangular grid} (Figure~\ref{figure:1.3.1}). Through each
point in the grid we draw a short line segment with slope
$f(x_i,y_j)$. The result is an approximation to a direction field  for
\eqref{eq:1.3.1} in $R$. If the grid points are sufficiently numerous and
close together, we can draw approximate integral curves of \eqref{eq:1.3.1}
by drawing curves through points in the grid tangent to the line
segments associated with the points in the grid.

\begin{figure}[H]
  \centering
  \includegraphics[width = 0.7\textwidth]{fig010301}
\color{blue}
  \caption{\quad  A rectangular grid}
  \label{figure:1.3.1}
\end{figure}

Unfortunately, approximating a direction field and graphing integral
curves in this way is too tedious to be done effectively by hand.
However, there is software  for doing this.
As you'll  see, the combination of  direction fields and
 integral curves gives  useful insights into the behavior of
the solutions of the differential equation even if
we can't obtain exact solutions.

We'll  study numerical methods for solving a single first order
equation \eqref{eq:1.3.1} in Chapter~3. These methods can be
used
to plot solution curves of \eqref{eq:1.3.1} in a rectangular region $R$
\emph{\color{blue}if $f$ is continuous on} $R$. Figures~\ref{figure:1.3.2},
\ref{figure:1.3.3}, and \ref{figure:1.3.4} show direction fields and solution
curves for the differential equations
\[
 y'=\frac{x^2-y^2}{1+x^2+y^2}, \quad
y'=1+xy^2,\text{\quad and \quad}
 y'=\frac{x-y}{1+x^2},
\]
which are all of the form \eqref{eq:1.3.1} with $f$
 continuous for all $(x,y)$.

\begin{figure}[htbp]
\color{blue}
  \begin{minipage}[b]{0.5\linewidth}
    \centering
   \scalebox{.6}{
  \includegraphics[%bb=-78 148 689 643,
  width=5.67in,height=3.66in,keepaspectratio]{fig010302}}
  \caption{\,A direction field and integral curves for
$y =\dst{\frac{x^{2}-y^{2}}{1+x^{2}+y^{2}}}$}
  \label{figure:1.3.2}
  \end{minipage}
  \hspace{0.6cm}
  \begin{minipage}[b]{0.5\linewidth}
    \centering
   \scalebox{.7}{
  \includegraphics[%bb=-78 148 689 643,
  width=5.67in,height=3.66in,keepaspectratio]{fig010303}}
\caption{A direction field and integral curves
for $y'=1+xy^2$}
  \label{figure:1.3.3}
  \end{minipage}
\end{figure}

\begin{figure}[tbp]
  \centering
  \includegraphics[%bb=-78 148 689 643,
  width=5.67in,height=3.66in,keepaspectratio]{fig010304}
\color{blue}
  \caption{A direction and integral curves for
$y'=\dst{\frac{x-y}{1+x^2}}$}
  \label{figure:1.3.4}
\end{figure}

The methods of Chapter~3 won't work for the equation
\begin{equation} \label{eq:1.3.2}
y'=-x/y
\end{equation}
if $R$ contains part of the $x$-axis, since $f(x,y)=-x/y$ is undefined
when $y=0$. Similarly, they won't work for the equation
\begin{equation} \label{eq:1.3.3}
y'=\frac{x^2}{1-x^2-y^2}
\end{equation}
if $R$ contains any part of the unit circle $x^2+y^2=1$, because the
right side of \eqref{eq:1.3.3} is undefined if $x^2+y^2=1$. However,
\eqref{eq:1.3.2} and \eqref{eq:1.3.3} can written as
\begin{equation} \label{eq:1.3.4}
y'=\frac{A(x,y)}{B(x,y)}
\end{equation}
where $A$ and $B$ are  continuous on any rectangle $R$. Because of
this,
some differential equation software is based on
numerically solving pairs of equations of the form
\begin{equation} \label{eq:1.3.5}
\frac{dx}{dt}=B(x,y),\quad \frac{dy}{dt}=A(x,y)
\end{equation}
where $x$ and $y$ are regarded as functions of a parameter $t$.
If $x=x(t)$ and $y=y(t)$  satisfy these equations, then
\[
y'=\frac{dy}{dx}=\frac{dy}{dt}\left/\frac{dx}{dt}\right.=\frac{A(x,y)}{B(x,y)},
\]
so $y=y(x)$ satisfies \eqref{eq:1.3.4}.



Eqns.~\eqref{eq:1.3.2} and \eqref{eq:1.3.3} can be reformulated as in
\eqref{eq:1.3.4} with
\[
\frac{dx}{dt}=-y,\quad \frac{dy}{dt}=x
\]
and
\[
\frac{dx}{dt}=1-x^2-y^2,\quad \frac{dy}{dt}=x^2,
\]
respectively. Even if $f$ is continuous and otherwise ``nice''
throughout $R$, your software  may require you to
reformulate the equation $y'=f(x,y)$ as
\[
\frac{dx}{dt}=1,\quad \frac{dy}{dt}=f(x,y),
\]
which is of the form \eqref{eq:1.3.5} with $A(x,y)=f(x,y)$ and
$B(x,y)=1$.

Figure~\ref{figure:1.3.5} shows a direction field and some integral curves
for \eqref{eq:1.3.2}. As we saw in
Example~\ref{example:1.2.1} and will verify
again in Section~2.2, the integral curves of \eqref{eq:1.3.2}
are circles centered at the origin.



\begin{figure}[H]
  \centering
  \includegraphics[%bb=-78 148 689 643,
  width=5.67in,height=3.66in,keepaspectratio]{fig010305}
\color{blue}
\caption{A direction field and integral curves for
$y'=-\dst\frac{x}{y}$}
  \label{figure:1.3.5}
\end{figure}

Figure~\ref{figure:1.3.6} shows a direction field and some integral curves
for \eqref{eq:1.3.3}. The integral curves near the top and bottom are
solution curves. However, the integral curves near the middle are more
complicated. For example, Figure~\ref{figure:1.3.7} shows the integral
curve through the origin. The vertices of the dashed rectangle are on
the circle $x^2+y^2=1$ ($a\approx.846$, $b\approx.533$), where all
integral curves of \eqref{eq:1.3.3} have infinite slope. There are
three solution curves of \eqref{eq:1.3.3} on the integral curve in the
figure: the segment above the level $y=b$ is the graph of a solution
on $(-\infty,a)$, the segment below the level $y=-b$ is the graph of a
solution on $(-a,\infty)$, and the segment between these two levels is
the graph of a solution on $(-a,a)$.

\technology
 \medskip
As you study from this book, you'll often be asked to use
computer software and graphics. Exercises with this intent are marked
as \Cex\, (computer or calculator required), \CGex\, (computer and/or
graphics required), or \Lex\ (laboratory work requiring software
and/or graphics).
 Often
you may not completely understand how the software does what it does.
This is similar to the situation most people are in when they drive
automobiles or watch television, and it doesn't decrease the value
of using modern technology as an aid to learning. Just be careful that
you use the technology as a supplement to thought rather than a
substitute for it.


\begin{figure}[htbp]
\color{blue}
  \begin{minipage}[b]{0.5\linewidth}
    \centering
   \scalebox{.7}{
  \includegraphics[%bb=-78 148 689 643,
  width=5.67in,height=3.66in,keepaspectratio]{fig010306} }
\caption{A direction field and integral curves for
$y'=\dst\frac{x^2}{-x^2-y^2}$}
  \label{figure:1.3.6}
  \end{minipage}
  \hspace{0.6cm}
  \begin{minipage}[b]{0.5\linewidth}
    \centering
   \scalebox{.7}{
  \includegraphics[%bb=-78 148 689 643,
  width=5.67in,height=3.66in,keepaspectratio]{fig010307} }
\caption{}
  \label{figure:1.3.7}
  \end{minipage}
\end{figure}

\enlargethispage{1in}

\exercises
\noindent
\emph{In Exercises~{\color{red}\rm\bf 1--11} a  direction
field is drawn for the given equation. Sketch some integral curves.}



\begin{figure}[H]
  \color{blue}
  \centering
  \includegraphics[%bb=-78 148 689 643,
  width=5.67in,height=3.66in,keepaspectratio]{exer010301}
  \caption*{{\color{red}\bf 1}\; A direction field for
$y'=\dst{\frac{x}{y}}$}
\end{figure}



\begin{figure}[H]
  \color{blue}
  \centering
  \includegraphics[%bb=-78 148 689 643,
  width=5.67in,height=3.66in,keepaspectratio]{exer010302}
  \caption*{{\color{red}\bf 2}\; A direction field for
 $\dst{y'=\dst\frac{2xy^2}{1+x^2}}$\quad}
\end{figure}






\begin{figure}[H]
  \color{blue}
  \centering
  \includegraphics[%bb=-78 148 689 643,
  width=5.67in,height=3.66in,keepaspectratio]{exer010303}
  \caption*{{\color{red}\bf 3}\; A direction field for
 $y'=x^2(1+y^2)$}
\end{figure}



\begin{figure}[H]
  \color{blue}
  \centering
  \includegraphics[%bb=-78 148 689 643,
  width=5.67in,height=3.66in,keepaspectratio]{exer010304}
  \caption*{{\color{red}\bf 4}\; A direction field for
$y'=\dst\frac{1}{1+x^2+y^2}$}
\end{figure}



\begin{figure}[H]
  \color{blue}
  \centering
  \includegraphics[%bb=-78 148 689 643,
  width=5.67in,height=3.66in,keepaspectratio]{exer010305}
  \caption*{{\color{red}\bf 5}\; A direction field for
$y'=-(2xy^2+y^3)$}
\end{figure}



\begin{figure}[H]
  \color{blue}
  \centering
  \includegraphics[%bb=-78 148 689 643,
  width=5.67in,height=3.66in,keepaspectratio]{exer010306}
  \caption*{{\color{red}\bf 6}\; A direction field for
$y'=(x^2+y^2)^{1/2}$}
\end{figure}



\begin{figure}[H]
  \color{blue}
  \centering
  \includegraphics[%bb=-78 148 689 643,
  width=5.67in,height=3.66in,keepaspectratio]{exer010307}
  \caption*{{\color{red}\bf 7}\; A direction field for
$y'=\sin xy$}
\end{figure}



\begin{figure}[H]
  \color{blue}
  \centering
  \includegraphics[%bb=-78 148 689 643,
  width=5.67in,height=3.66in,keepaspectratio]{exer010308}
  \caption*{{\color{red}\bf 8}\; A direction field for
 $y'=e^{xy}$}
\end{figure}



\begin{figure}[H]
  \color{blue}
  \centering
  \includegraphics[%bb=-78 148 689 643,
  width=5.67in,height=3.66in,keepaspectratio]{exer010309}
  \caption*{{\color{red}\bf 9}\; A direction field for
$y'=(x-y^2)(x^2-y)$}
\end{figure}



\begin{figure}[H]
  \color{blue}
  \centering
  \includegraphics[%bb=-78 148 689 643,
  width=5.67in,height=3.66in,keepaspectratio]{exer010310}
  \caption*{{\color{red}\bf 10}\; A direction field for
$y'=x^3y^2+xy^3$}
\end{figure}



\begin{figure}[H]
  \color{blue}
  \centering
  \includegraphics[%bb=-78 148 689 643,
  width=5.67in,height=3.66in,keepaspectratio]{exer010311}
  \caption*{{\color{red}\bf 11}\; A direction field for
$y'=\sin(x-2y)$}
\end{figure}

\emph{In Exercises~\ref{exer:1.3.12}-\ref{exer:1.3.22} construct a
direction
field and plot some integral curves in the indicated rectangular region.}


\begin{exerciselist}
\setcounter{exercise}{11}
\item\label{exer:1.3.12}\CGex
$y'=y(y-1);    \quad  \{-1\le x\le 2,\ -2\le y\le2\}$

\item\label{exer:1.3.13} \CGex
$y'=2-3xy;    \quad  \{-1\le x\le 4,\ -4\le y\le4\}$

\item\label{exer:1.3.14} \CGex
$y'=xy(y-1);    \quad  \{-2\le x\le2,\ -4\le y\le 4\}$

\item\label{exer:1.3.15}\CGex
$y'=3x+y;    \quad  \{-2\le x\le2,\ 0\le y\le 4\}$

\item\label{exer:1.3.16} \CGex
$y'=y-x^3;    \quad  \{-2\le x\le2,\ -2\le y\le 2\}$

\item\label{exer:1.3.17}\CGex
$y'=1-x^2-y^2;    \quad  \{-2\le x\le2,\ -2\le y\le 2\}$

\item\label{exer:1.3.18}\CGex
$y'=x(y^2-1);    \quad  \{-3\le x\le3,\ -3\le y\le 2\}$

\item\label{exer:1.3.19}\CGex
$y'=\dst\frac{x}{y(y^2-1)};    \quad  \{-2\le x\le2,\ -2\le y\le 2\}$

\item\label{exer:1.3.20}\CGex
$y'=\dst\frac{xy^2}{y-1};    \quad  \{-2\le x\le2,\ -1\le y\le 4\}$

\item\label{exer:1.3.21}\CGex
$y'=\dst\frac{x(y^2-1)}{y};    \quad  \{-1\le x\le1,\ -2\le y\le 2\}$

\item\label{exer:1.3.22}\CGex
$y'=-\dst\frac{x^2+y^2}{1-x^2-y^2};    \quad  \{-2\le x\le2,\ -2\le y\le 2\}$


\item\label{exer:1.3.23} \Lex
By suitably renaming the constants and  dependent variables
in the equations
\[
T' = -k(T-T_m)
\eqno{\rm(A)}
\]
and
\[
G'=-\lambda G+r
\eqno{\rm(B)}
\]
discussed in Section~1.2 in connection with Newton's
law of cooling and absorption of glucose in the body, we can
write both as
\[
y'=- ay+b,
\eqno{\rm(C)}
\]
where $a$ is a positive constant and $b$ is an arbitrary
constant. Thus, (A) is of the form (C) with $y=T$, $a=k$, and
$b=kT_m$, and (B) is of the form (C) with $y=G$,
$a=\lambda$, and $b=r$. We'll encounter equations of the form
(C) in many other applications in Chapter~2.

Choose a positive  $a$  and an
arbitrary  $b$.
 Construct a  direction field and plot some integral curves for (C) in
a rectangular
region of the form
\[
\{0\le t\le T,\ c\le y\le d\}
\]
 of the $ty$-plane.
 Vary $T$, $c$, and $d$ until you discover a common property
of all the solutions of (C). Repeat this experiment with various
choices of $a$ and $b$ until you can state
this property precisely in terms of $a$ and $b$.

\item\label{exer:1.3.24} \Lex
By suitably renaming the constants and  dependent variables
in the equations
\[
P'=aP(1-\alpha P)
\eqno{\rm(A)}
\]
and
\[
I'=rI(S-I)
\eqno{\rm(B)}
\]
discussed in Section~1.1 in connection with Verhulst's
population
model and the spread of an epidemic, we can write both in the form
\[
y'=ay-by^2,
\eqno{\rm(C)}
\]
where $a$ and $b$ are positive constants. Thus, (A) is of the form (C)
with $y=P$, $a=a$, and $b=a\alpha$, and (B) is of the form (C) with
$y=I$, $a=rS$, and $b=r$. In Chapter~2 we'll encounter
equations of the form
(C) in
 other applications..

\begin{alist}
\item % (a)
Choose positive numbers $a$ and $b$. Construct a direction field and
plot some integral curves for (C) in a rectangular region of the form
\[
\{0\le t\le T,\ 0\le y\le d\}
\]
of the $ty$-plane. Vary $T$ and $d$ until you discover a common
property of all solutions of (C) with $y(0)>0$. Repeat this experiment
with various choices of $a$ and $b$ until you can state this property
precisely in terms of $a$ and $b$.

\item % (a)
Choose  positive numbers $a$ and $b$. Construct a direction field and
plot some integral curves for (C) in a rectangular region of the form
\[
\{0\le t\le T,\ c\le y\le 0\}
\]
of the $ty$-plane. Vary $a$, $b$, $T$ and $c$ until you discover a
common property of all solutions of (C) with $y(0)<0$.

You can verify your results later by doing Exercise~2.2.~\hspace*{-3pt}\ref{exer:2.2.27}.

\end{alist}
\end{exerciselist}
